%%%%%%%%%%%%%%%%%%%%% chapter.tex %%%%%%%%%%%%%%%%%%%%%%%%%%%%%%%%%
%
% sample chapter
%
% Use this file as a template for your own input.
%
%%%%%%%%%%%%%%%%%%%%%%%% Springer-Verlag %%%%%%%%%%%%%%%%%%%%%%%%%%
%\motto{Use the template \emph{chapter.tex} to style the various elements of your chapter content.}

\chapter{Rosetta Code Tasks starting with B}

\section*{Balanced brackets}


\textbf{Task}:

\begin{itemize}
\item
  Generate a string with N opening brackets (``\texttt{{[}}'') and N
  closing brackets (``\texttt{{]}}''), in some arbitrary order.
\item
  Determine whether the generated string is \emph{balanced}; that is,
  whether it consists entirely of pairs of opening/closing brackets (in
  that order), none of which mis-nest.
\end{itemize}

\textbf{Examples}:

\begin{verbatim}
   (empty)   OK
   []        OK   ][        NOT OK
   [][]      OK   ][][      NOT OK
   [[][]]    OK   []][[]    NOT OK
\end{verbatim}



\begin{wideverbatim}

(load "@lib/simul.l")  # For 'shuffle'

(de generateBrackets (N)
   (shuffle (make (do N (link "[" "]")))) )

(de checkBrackets (S)
   (let N 0
      (for C S
         (if (= C "[")
            (inc 'N)
            (if2 (= C "]") (=0 N)
               (off N)
               (dec 'N) ) ) )
      (=0 N) ) )

(for N 10
   (prinl (if (checkBrackets (prin (generateBrackets N))) " OK" "not OK")) )

Output:

[] OK
[[]] OK
]]][[[not OK
[[[][]]] OK
[][][[[]]] OK
[]][[[][[]]]not OK
[[[]]][][][][] OK
]][][[[[]][]]][[not OK
[]][][[[][[]]][]][not OK
[[[][]]]]][][[]]][[[not OK

\end{wideverbatim}

\pagebreak{}
\section*{Best shuffle}


Shuffle the characters of a string in such a way that as many of the
character values are in a different position as possible. Print the
result as follows: original string, shuffled string, (score). The score
gives the number of positions whose character value did \emph{not}
change.

For example: \texttt{tree, eetr, (0)}

A shuffle that produces a randomized result among the best choices is to
be preferred. A deterministic approach that produces the same sequence
every time is acceptable as an alternative.

The words to test with are: \texttt{abracadabra}, \texttt{seesaw},
\texttt{elk}, \texttt{grrrrrr}, \texttt{up}, \texttt{a}

Cf.

\begin{itemize}
\item
  \emph{Anagrams/Deranged anagrams}
\item
  \emph{Permutations/Derangements}
\end{itemize}


\begin{wideverbatim}

(de bestShuffle (Str)
   (let Lst NIL
      (for C (setq Str (chop Str))
         (if (assoc C Lst)
            (con @ (cons C (cdr @)))
            (push 'Lst (cons C)) ) )
      (setq Lst (apply conc (flip (by length sort Lst))))
      (let Res
         (mapcar
            '((C)
               (prog1 (or (find <> Lst (circ C)) C)
                  (setq Lst (delete @ Lst)) ) )
            Str )
         (prinl Str " " Res " (" (cnt = Str Res) ")") ) ) )

Output:

: (bestShuffle "abracadabra")
abracadabra raarababadc (0)

: (bestShuffle "seesaw")
seesaw essewa (0)

: (bestShuffle "elk")
elk lke (0)

: (bestShuffle "grrrrrr")
grrrrrr rgrrrrr (5)

: (bestShuffle "up")
up pu (0)

: (bestShuffle "a")
a a (1)

\end{wideverbatim}

\pagebreak{}
\section*{Binary digits}


The task is to output the sequence of binary digits for a given
\href{http://en.wikipedia.org/wiki/Natural\_number}{non-negative
integer}.

The decimal value \texttt{5}, should produce an output of \texttt{101}
The decimal value \texttt{50} should produce an output of
\texttt{110010} The decimal value \texttt{9000} should produce an output
of \texttt{10001100101000}

The results can be achieved using builtin radix functions within the
language, if these are available, or alternatively a user defined
function can be used. The output produced should consist just of the
binary digits of each number followed by a newline. There should be no
other whitespace, radix or sign markers in the produced output, and
\href{http://en.wikipedia.org/wiki/Leading\_zero}{leading zeros} should
not appear in the results.



\begin{wideverbatim}

: (bin 5)
-> "101"

: (bin 50)
-> "110010"

: (bin 9000)
-> "10001100101000"

\end{wideverbatim}

\pagebreak{}
\section*{Binary search}


A binary search divides a range of values into halves, and continues to
narrow down the field of search until the unknown value is found. It is
the classic example of a ``divide and conquer'' algorithm.

As an analogy, consider the children's game "\emph{guess a number}."
The scorer has a secret number, and will only tell the player if their
guessed number is higher than, lower than, or equal to the secret
number. The player then uses this information to guess a new number.

As the player, an optimal strategy for the general case is to start by
choosing the range's midpoint as the guess, and then asking whether the
guess was higher, lower, or equal to the secret number. If the guess was
too high, one would select the point exactly between the range midpoint
and the beginning of the range. If the original guess was too low, one
would ask about the point exactly between the range midpoint and the end
of the range. This process repeats until one has reached the secret
number.

\textbf{The Task}

Given the starting point of a range, the ending point of a range, and
the ``secret value'', implement a binary search through a sorted integer
array for a certain number. Implementations can be recursive or
iterative (both if you can). Print out whether or not the number was in
the array afterwards. If it was, print the index also.

There are several binary search algorithms commonly seen. They differ by
how they treat multiple values equal to the given value, and whether
they indicate whether the element was found or not. For completeness we
will present pseudocode for all of them.

All of the following code examples use an ``inclusive'' upper bound
(i.e. \texttt{high = N-1} initially). Any of the examples can be
converted into an equivalent example using ``exclusive'' upper bound
(i.e. \texttt{high = N} initially) by making the following simple
changes (which simply increase \texttt{high} by 1):

\begin{itemize}
\item
  change \texttt{high = N-1} to \texttt{high = N}
\item
  change \texttt{high = mid-1} to \texttt{high = mid}
\item
  (for recursive algorithm) change \texttt{if (high \textless{} low)} to
  \texttt{if (high \textless{}= low)}
\item
  (for iterative algorithm) change
  \texttt{while (low \textless{}= high)} to
  \texttt{while (low \textless{} high)}
\end{itemize}

Traditional algorithm

The algorithms are as follows (from
\href{http://en.wikipedia.org/wiki/Binary\_search}{Wikipedia}). The
algorithms return the index of some element that equals the given value
(if there are multiple such elements, it returns some arbitrary one). It
is also possible, when the element is not found, to return the
``insertion point'' for it (the index that the value would have if it
were inserted into the array).

\textbf{Recursive Pseudocode}:

\begin{wideverbatim}
  // initially called with low = 0, high = N-1
  BinarySearch(A[0..N-1], value, low, high) {
      // invariants: value > A[i] for all i < low
                     value < A[i] for all i > high
      if (high < low)
          return not_found // value would be inserted at index "low"
      mid = (low + high) / 2
      if (A[mid] > value)
          return BinarySearch(A, value, low, mid-1)
      else if (A[mid] < value)
          return BinarySearch(A, value, mid+1, high)
      else
          return mid
  }
\end{wideverbatim}

\textbf{Iterative Pseudocode}:

\begin{wideverbatim}
  BinarySearch(A[0..N-1], value) {
      low = 0
      high = N - 1
      while (low <= high) {
          // invariants: value > A[i] for all i < low
                         value < A[i] for all i > high
          mid = (low + high) / 2
          if (A[mid] > value)
              high = mid - 1
          else if (A[mid] < value)
              low = mid + 1
          else
              return mid
      }
      return not_found // value would be inserted at index "low"
  }
\end{wideverbatim}

Leftmost insertion point

The following algorithms return the leftmost place where the given
element can be correctly inserted (and still maintain the sorted order).
This is the lower (inclusive) bound of the range of elements that are
equal to the given value (if any). Equivalently, this is the lowest
index where the element is greater than or equal to the given value
(since if it were any lower, it would violate the ordering), or 1 past
the last index if such an element does not exist. This algorithm does
not determine if the element is actually found. This algorithm only
requires one comparison per level.

\textbf{Recursive Pseudocode}:

\begin{wideverbatim}
  // initially called with low = 0, high = N - 1
  BinarySearch_Left(A[0..N-1], value, low, high) {
      // invariants: value > A[i] for all i < low
                     value <= A[i] for all i > high
      if (high < low)
          return low
      mid = (low + high) / 2
      if (A[mid] >= value)
          return BinarySearch_Left(A, value, low, mid-1)
      else
          return BinarySearch_Left(A, value, mid+1, high)
  }
\end{wideverbatim}

\textbf{Iterative Pseudocode}:

\begin{wideverbatim}
  BinarySearch_Left(A[0..N-1], value) {
      low = 0
      high = N - 1
      while (low <= high) {
          // invariants: value > A[i] for all i < low
                         value <= A[i] for all i > high
          mid = (low + high) / 2
          if (A[mid] >= value)
              high = mid - 1
          else
              low = mid + 1
      }
      return low
  }
\end{wideverbatim}

Rightmost insertion point

The following algorithms return the rightmost place where the given
element can be correctly inserted (and still maintain the sorted order).
This is the upper (exclusive) bound of the range of elements that are
equal to the given value (if any). Equivalently, this is the lowest
index where the element is greater than the given value, or 1 past the
last index if such an element does not exist. This algorithm does not
determine if the element is actually found. This algorithm only requires
one comparison per level. Note that these algorithms are almost exactly
the same as the leftmost-insertion-point algorithms, except for how the
inequality treats equal values.

\textbf{Recursive Pseudocode}:

\begin{wideverbatim}
  // initially called with low = 0, high = N - 1
  BinarySearch_Right(A[0..N-1], value, low, high) {
      // invariants: value >= A[i] for all i < low
                     value < A[i] for all i > high
      if (high < low)
          return low
      mid = (low + high) / 2
      if (A[mid] > value)
          return BinarySearch_Right(A, value, low, mid-1)
      else
          return BinarySearch_Right(A, value, mid+1, high)
  }
\end{wideverbatim}

\textbf{Iterative Pseudocode}:

\begin{wideverbatim}
  BinarySearch_Right(A[0..N-1], value) {
      low = 0
      high = N - 1
      while (low <= high) {
          // invariants: value >= A[i] for all i < low
                         value < A[i] for all i > high
          mid = (low + high) / 2
          if (A[mid] > value)
              high = mid - 1
          else
              low = mid + 1
      }
      return low
  }
\end{wideverbatim}

Extra credit

Make sure it does not have overflow bugs.

The line in the pseudocode above to calculate the mean of two integers:

\begin{wideverbatim}
mid = (low + high) / 2
\end{wideverbatim}

could produce the wrong result in some programming languages when used
with a bounded integer type, if the addition causes an overflow. (This
can occur if the array size is greater than half the maximum integer
value.) If signed integers are used, and \texttt{low + high} overflows,
it becomes a negative number, and dividing by 2 will still result in a
negative number. Indexing an array with a negative number could produce
an out-of-bounds exception, or other undefined behavior. If unsigned
integers are used, an overflow will result in losing the largest bit,
which will produce the wrong result.

One way to fix it is to manually add half the range to the low number:

\begin{wideverbatim}
mid = low + (high - low) / 2
\end{wideverbatim}

Even though this is mathematically equivalent to the above, it is not
susceptible to overflow.

Another way for signed integers, possibly faster, is the following:

\begin{wideverbatim}
mid = (low + high) >>> 1
\end{wideverbatim}

where \texttt{ \textgreater{}\textgreater{}\textgreater{} } is the
logical right shift operator. The reason why this works is that, for
signed integers, even though it overflows, when viewed as an unsigned
number, the value is still the correct sum. To divide an unsigned number
by 2, simply do a logical right shift.

\textbf{References:}\\

\begin{itemize}
\item
  C.f: \emph{Guess the number/With Feedback (Player)}
\item
  \href{http://en.wikipedia.org/wiki/Binary\_search\_algorithm}{wp:Binary
  search algorithm}
\item
  \href{http://googleresearch.blogspot.com/2006/06/extra-extra-read-all-about-it-nearly.html}{Extra,
  Extra - Read All About It: Nearly All Binary Searches and Mergesorts
  are Broken}.
\end{itemize}



\begin{wideverbatim}

(de recursiveSearch (Val Lst Len)
   (unless (=0 Len)
      (let (N (inc (/ Len 2))  L (nth Lst N))
         (cond
            ((= Val (car L)) Val)
            ((> Val (car L))
               (recursiveSearch Val (cdr L) (- Len N)) )
            (T (recursiveSearch Val Lst (dec N))) ) ) ) )

Output:

: (recursiveSearch 5 (2 3 5 8 "abc" "klm" "xyz" (7) (a b)) 9)
-> 5
: (recursiveSearch '(a b) (2 3 5 8 "abc" "klm" "xyz" (7) (a b)) 9)
-> (a b)
: (recursiveSearch (9) (2 3 5 8 "abc" "klm" "xyz" (7) (a b)) 9)
-> NIL

(de iterativeSearch (Val Lst Len)
   (use (N L)
      (loop
         (T (=0 Len))
         (setq
            N (inc (/ Len 2))
            L (nth Lst N) )
         (T (= Val (car L)) Val)
         (if (> Val (car L))
            (setq Lst (cdr L)  Len (- Len N))
            (setq Len (dec N)) ) ) ) )

Output:

: (iterativeSearch 5 (2 3 5 8 "abc" "klm" "xyz" (7) (a b)) 9)
-> 5
: (iterativeSearch '(a b) (2 3 5 8 "abc" "klm" "xyz" (7) (a b)) 9)
-> (a b)
: (iterativeSearch (9) (2 3 5 8 "abc" "klm" "xyz" (7) (a b)) 9)
-> NIL

\end{wideverbatim}

\pagebreak{}
\section*{Binary strings}


Many languages have powerful and useful (\textbf{binary safe})
\href{http://en.wikipedia.org/wiki/String\_(computer\_science)}{string}
\href{http://en.wikipedia.org/wiki/Comparison\_of\_programming\_languages\_(string\_functions)}{manipulation
functions}, while others don't, making it harder for these languages to
accomplish some tasks. This task is about creating functions to handle
\emph{binary} strings (strings made of arbitrary bytes, i.e. \emph{byte
strings} according to Wikipedia) for those languages that don't have
built-in support for them. If your language of choice does have this
built-in support, show a possible alternative implementation for the
\emph{functions} or \emph{abilities} already provided by the language.
In particular the functions you need to create are:

\begin{itemize}
\item
  String creation and destruction (when needed and if there's no
  \emph{garbage collection} or similar
  mechanism)
\item
  String assignment
\item
  String comparison
\item
  String cloning and copying
\item
  Check if a string is empty
\item
  Append a byte to a string
\item
  Extract a substring from a string
\item
  Replace every occurrence of a byte (or a string) in a string with
  another string
\item
  Join strings
\end{itemize}

Possible contexts of use: compression algorithms (like \emph{LZW
  compression}), L-systems (manipulation of symbols), many more.


\begin{wideverbatim}

Byte strings are represented in PicoLisp as lists of numbers. They can be
maniplated easily with the built-in list functionality.

I/O of raw bytes is done via the 'wr' (write) and 'rd' (read) functions. The
following creates a file consisting of 256 bytes, with values from 0 to 255:

: (out "rawfile"
   (mapc wr (range 0 255)) )

Looking at a hex dump of that file:

: (hd "rawfile")
00000000  00 01 02 03 04 05 06 07 08 09 0A 0B 0C 0D 0E 0F  ................
00000010  10 11 12 13 14 15 16 17 18 19 1A 1B 1C 1D 1E 1F  ................
00000020  20 21 22 23 24 25 26 27 28 29 2A 2B 2C 2D 2E 2F   !"#\$\%\&'()*+,-./
00000030  30 31 32 33 34 35 36 37 38 39 3A 3B 3C 3D 3E 3F  0123456789:;<=>?
...


To read part of that file, an external tool like 'dd' might be used:

: (in '(dd "skip=32" "bs=1" "count=16" "if=rawfile")
   (make
      (while (rd 1)
         (link @) ) ) )
-> (32 33 34 35 36 37 38 39 40 41 42 43 44 45 46 47)

Now such byte lists can be assigned the normal way ('let', 'setq' etc.), they
can be compared with '=', '>', '>=' etc, and manipulated with all internal map-,
filter-, concatenation-, reversal-, pattern matching, and other functions.

If desired, a string containing meaningful values can also be converted to
a transient symbol, e.g. the example above

: (pack (mapcar char (32 33 34 35 36 37 38 39 40 41 42 43 44 45 46 47)))
-> " !\"#\$\%\&'()*+,-./"

\end{wideverbatim}

\pagebreak{}
\section*{Bitmap}


Show a basic storage type to handle a simple RGB raster graphics image,
and some primitive associated functions.

If possible provide a function to allocate an uninitialised image, given
its width and height, and provide 3 additional functions:

\begin{itemize}
\item
  one to fill an image with a plain RGB color,
\item
  one to set a given pixel with a color,
\item
  one to get the color of a pixel.
\end{itemize}

(If there are specificities about the storage or the allocation, explain
those.)

\emph{These functions are used as a base for the articles in the
  category \emph{raster graphics operations}, and a basic output
  function to check the results is available in the article
  \emph{write ppm file}.}

\begin{wideverbatim}

For time critical applications this would be done with inline-C in PicoLisp,
but especially for small bitmaps the following makes sense.

# Create an empty image of 120 x 90 pixels
(setq *Ppm (make (do 90 (link (need 120)))))

# Fill an image with a given color
(de ppmFill (Ppm R G B)
   (for Y Ppm
      (map
         '((X) (set X (list R G B)))
         Y ) ) )

# Set pixel with a color
(de ppmSetPixel (Ppm X Y R G B)
   (set (nth Ppm Y X) (list R G B)) )

# Get the color of a pixel
(de ppmGetPixel (Ppm X Y)
   (get Ppm Y X) )

\end{wideverbatim}

\pagebreak{}
\section*{Bitmap/B\'{e}zier curves/Cubic}


Using the data storage type defined in \emph{Basic bitmap storage}
for raster images, and the \texttt{draw\_line} function defined in
\emph{Bresenham's line algorithm}, draw a \textbf{cubic bezier
  curve}
(\href{http://en.wikipedia.org/wiki/Bezier\_curves\#Cubic\_B.C3.A9zier\_curves}{definition
  on Wikipedia}).

\begin{wideverbatim}

This uses the 'brez' line drawing function from
[[Bitmap/Bresenham's line algorithm#PicoLisp]].

(scl 6)

(de cubicBezier (Img N X1 Y1 X2 Y2 X3 Y3 X4 Y4)
   (let (R (* N N N)  X X1  Y Y1  DX 0  DY 0)
      (for I N
         (let
            (J (- N I)
               A (*/ 1.0 J J J R)
               B (*/ 3.0 I J J R)
               C (*/ 3.0 I I J R)
               D (*/ 1.0 I I I R) )
            (brez Img
               X
               Y
               (setq DX
                  (-
                     (+ (*/ A X1 1.0) (*/ B X2 1.0) (*/ C X3 1.0) (*/ D X4 1.0))
                     X ) )
               (setq DY
                  (-
                     (+ (*/ A Y1 1.0) (*/ B Y2 1.0) (*/ C Y3 1.0) (*/ D Y4 1.0))
                     Y ) ) )
            (inc 'X DX)
            (inc 'Y DY) ) ) ) )

\end{wideverbatim}

\begin{wideverbatim}

Test:

(let Img (make (do 200 (link (need 300 0))))       # Create image 300 x 200
   (cubicBezier Img 24 20 120 540 33 -225 33 285 100)
   (out "img.pbm"                                  # Write to bitmap file
      (prinl "P1")
      (prinl 300 " " 200)
      (mapc prinl Img) ) )

(call 'display "img.pbm")

\end{wideverbatim}

\pagebreak{}
\section*{Bitmap/B\'{e}zier curves/Quadratic}

Using the data storage type defined in \emph{Basic bitmap storage}
for raster images, and the \texttt{draw\_line} function defined in
\emph{Bresenham's line algorithm}, draw a \textbf{quadratic bezier
  curve} 
(\href{http://en.wikipedia.org/wiki/Bezier\_curves\#Quadratic\_B.C3.A9zier\_curves}{definition
on Wikipedia}).

\begin{wideverbatim}

This uses the 'brez' line drawing function from
[[Bitmap/Bresenham's line algorithm#PicoLisp]].

(scl 6)

(de quadBezier (Img N X1 Y1 X2 Y2 X3 Y3)
   (let (R (* N N)  X X1  Y Y1  DX 0  DY 0)
      (for I N
         (let (J (- N I)  A (*/ 1.0 J J R)  B (*/ 2.0 I J R)  C (*/ 1.0 I I R))
            (brez Img X Y
               (setq DX (- (+ (*/ A X1 1.0) (*/ B X2 1.0) (*/ C X3 1.0)) X))
               (setq DY (- (+ (*/ A Y1 1.0) (*/ B Y2 1.0) (*/ C Y3 1.0)) Y)) )
            (inc 'X DX)
            (inc 'Y DY) ) ) ) )

Test:

(let Img (make (do 200 (link (need 300 0))))       # Create image 300 x 200
   (quadBezier Img 12 20 100 300 -80 260 180)
   (out "img.pbm"                                  # Write to bitmap file
      (prinl "P1")
      (prinl 300 " " 200)
      (mapc prinl Img) ) )

(call 'display "img.pbm")

\end{wideverbatim}

\pagebreak{}
\section*{Bitmap/Bresenham's line algorithm}

Using the data storage type defined in \emph{Basic bitmap storage}
for raster images, draw a line given 2 points with the
\href{http://en.wikipedia.org/wiki/Bresenham\%27s\_line\_algorithm}{Bresenham's
algorithm}.

\begin{wideverbatim}

(de brez (Img X Y DX DY)
   (let SX
      (cond
         ((=0 DX) 0)
         ((gt0 DX) 1)
         (T (setq DX (- DX)) -1) )
      (let SY
         (cond
            ((=0 DY) 0)
            ((gt0 DY) 1)
            (T (setq DY (- DY)) -1) )
         (if (>= DX DY)
            (let E (- (* 2 DY) DX)
               (do DX
                  (set (nth Img Y X) 1)
                  (when (ge0 E)
                     (inc 'Y SY)
                     (dec 'E (* 2 DX)) )
                  (inc 'X SX)
                  (inc 'E (* 2 DY)) ) )
            (let E (- (* 2 DX) DY)
               (do DY
                  (set (nth Img Y X) 1)
                  (when (ge0 E)
                     (inc 'X SX)
                     (dec 'E (* 2 DY)) )
                  (inc 'Y SY)
                  (inc 'E (* 2 DX)) ) ) ) ) ) )

(let Img (make (do 90 (link (need 120 0))))        # Create image 120 x 90
   (brez Img 10 10 100 30)                         # Draw five lines
   (brez Img 10 10 100 50)
   (brez Img 10 10 100 70)
   (brez Img 10 10 60 70)
   (brez Img 10 10 20 70)
   (out "img.pbm"                                  # Write to bitmap file
      (prinl "P1")
      (prinl 120 " " 90)
      (mapc prinl Img) ) )

\end{wideverbatim}

\pagebreak{}
\section*{Bitmap/Flood fill}


Implement a \href{http://en.wikipedia.org/wiki/flood\_fill}{flood fill}.

A flood fill is a way of filling an area using \emph{color banks} to
define the contained area or a \emph{target color} which ``determines''
the area (the \emph{valley} that can be flooded; Wikipedia uses the term
\emph{target color}). It works almost like a water flooding from a point
towards the banks (or: inside the valley): if there's a hole in the
banks, the flood is not contained and all the image (or all the
``connected valleys'') get filled.

To accomplish the task, you need implementing just one of the possible
algorithms (examples are on
\href{http://en.wikipedia.org/wiki/flood\_fill}{Wikipedia}). Variations
on the \emph{theme} are allowed (e.g. adding a tolerance parameter or
argument for color-matching of the \emph{banks} or \emph{target} color).

{\includegraphics{graphics/128px-Unfilledcirc.png}}

\textbf{Testing}: the basic algorithm is not suitable for
\emph{truecolor} images; a possible test image is the one shown on the
right box; you can try to fill the white area, or the black inner
circle.

\begin{wideverbatim}

Using the format of [[Bitmap#PicoLisp|Bitmap]], a minimal recursive solution:

(de ppmFloodFill (Ppm X Y Color)
   (let Target (get Ppm Y X)
      (recur (X Y)
         (when (= Target (get Ppm Y X))
            (set (nth Ppm Y X) Color)
            (recurse (dec X) Y)
            (recurse (inc X) Y)
            (recurse X (dec Y))
            (recurse X (inc Y)) ) ) )
   Ppm )

Test using 'ppmRead' from [[Bitmap/Read a PPM file#PicoLisp]] and
'ppmWrite' from [[Bitmap/Write a PPM file#PicoLisp]], 
filling the white area with red:

(ppmWrite
   (ppmFloodFill (ppmRead "Unfilledcirc.ppm") 192 128 (255 0 0))
   "Filledcirc.ppm" )

\end{wideverbatim}

\pagebreak{}
\section*{Bitmap/Histogram}


Extend the basic bitmap storage defined \emph{on this page} to support
dealing with image histograms. The image histogram contains for each
luminance the count of image pixels having this luminance. Choosing a
histogram representation take care about the data type used for the
counts. It must have range of at least 0..NxM, where N is the image
width and M is the image height.

\textbf{Test task}

Histogram is useful for many image processing operations. As an example,
use it to convert an image into black and white art. The method works as
follows:

\begin{itemize}
\item
  Convert image to grayscale;
\item
  Compute the histogram
\item
  Find the median: defined as the luminance such that the image has an
  approximately equal number of pixels with lesser and greater
  luminance.
\item
  Replace each pixel of luminance lesser than the median to black, and
  others to white.
\end{itemize}

Use \emph{read}/\emph{write ppm file}, and \emph{grayscale image}
solutions.


\begin{wideverbatim}



(de histogram (Pgm)
   (let H (need 256 0)
      (for L Pgm
         (for G L
            (inc (nth H (inc G))) ) )
      H ) )

\end{wideverbatim}

\pagebreak{}
\section*{Bitmap/Midpoint circle algorithm}

Using the data storage type defined \emph{Basic bitmap storage} for
raster images, write an implementation of the \textbf{midpoint circle
  algorithm} (also known as \textbf{Bresenham's circle algorithm}).
(\href{http://en.wikipedia.org/wiki/Midpoint\_circle\_algorithm}{definition
  on Wikipedia}).


\begin{wideverbatim}

(de midPtCircle (Img CX CY Rad)
   (let (F (- 1 Rad)  DdFx 0  DdFy (* -2 Rad)  X 0  Y Rad)
      (set (nth Img (+ CY Rad) CX) 1)
      (set (nth Img (- CY Rad) CX) 1)
      (set (nth Img CY (+ CX Rad)) 1)
      (set (nth Img CY (- CX Rad)) 1)
      (while (> Y X)
         (when (ge0 F)
            (dec 'Y)
            (inc 'F (inc 'DdFy 2)) )
         (inc 'X)
         (inc 'F (inc (inc 'DdFx 2)))
         (set (nth Img (+ CY Y) (+ CX X)) 1)
         (set (nth Img (+ CY Y) (- CX X)) 1)
         (set (nth Img (- CY Y) (+ CX X)) 1)
         (set (nth Img (- CY Y) (- CX X)) 1)
         (set (nth Img (+ CY X) (+ CX Y)) 1)
         (set (nth Img (+ CY X) (- CX Y)) 1)
         (set (nth Img (- CY X) (+ CX Y)) 1)
         (set (nth Img (- CY X) (- CX Y)) 1) ) ) )

(let Img (make (do 120 (link (need 120 0))))       # Create image 120 x 120
   (midPtCircle Img 60 60 50)                      # Draw circle
   (out "img.pbm"                                  # Write to bitmap file
      (prinl "P1")
      (prinl 120 " " 120)
      (mapc prinl Img) ) )

\end{wideverbatim}

\pagebreak{}
\section*{Bitmap/PPM conversion through a pipe}


Using the data storage type defined \emph{Basic bitmap storage} for
for raster images, delegate writing a JPEG file through a
\textbf{pipe} using the \texttt{output\_ppm} function defined
\emph{Write ppm file}.

There are various utilities that can be used for this task, for example:
\textbf{cjpeg} (package \emph{``jpeg-progs''} on Linux),
\textbf{ppmtojpeg} (package \emph{``netpbm''} on Linux),
\textbf{convert} (from \emph{ImageMagick}, multi-platform).


\begin{wideverbatim}

# Create an empty image of 120 x 90 pixels
(setq *Ppm (make (do 90 (link (need 120)))))

# Fill background with green color
(ppmFill *Ppm 0 255 0)

# Draw a diagonal line
(for I 80 (ppmSetPixel *Ppm I I 0 0 0))

# Write to "img.jpg" through a pipe
(ppmWrite *Ppm '("convert" "-" "img.jpg"))

\end{wideverbatim}

\pagebreak{}
\section*{Bitmap/Read a PPM file}

Using the data storage type defined \emph{Basic bitmap storage} for
raster images, read an image from a PPM file (binary P6 prefered).
(Read \href{http://en.wikipedia.org/wiki/Netpbm\_format}{the
  definition of PPM file} on Wikipedia.)

\textbf{Task}: Use \emph{write ppm file} solution and \emph{grayscale
  image} solution with this one in order to convert a color image to
grayscale one.

\begin{wideverbatim}

(de ppmRead (File)
   (in File
      (unless (and `(hex "5036") (rd 2))  # P6
         (quit "Wrong file format" File) )
      (rd 1)
      (let (DX 0  DY 0  Max 0  C)
         (while (>= 9 (setq C (- (rd 1) `(char "0"))) 0)
            (setq DX (+ (* 10 DX) C)) )
         (while (>= 9 (setq C (- (rd 1) `(char "0"))) 0)
            (setq DY (+ (* 10 DY) C)) )
         (while (>= 9 (setq C (- (rd 1) `(char "0"))) 0)
            (setq Max (+ (* 10 Max) C)) )
         (prog1
            (make (do DY (link (need DX))))
            (for Y @
               (map
                  '((X) (set X (list (rd 1) (rd 1) (rd 1))))
                  Y ) ) ) ) ) )

Read a color image "img.ppm", convert and write to "img.pgm":

(pgmWrite (ppm->pgm (ppmRead "img.ppm")) "img.pgm")

\end{wideverbatim}

\pagebreak{}
\section*{Bitmap/Read an image through a pipe}


This task is the \emph{opposite} of the \emph{PPM conversion through a
  pipe}. In this task, using a delegate tool (like \textbf{cjpeg}, one
of the netpbm package, or \textbf{convert} of the ImageMagick package)
we read an image file and load it into the data storage type
\emph{Basic bitmap storage}. We can also use the code from \emph{Read
  ppm file}, so that we can use PPM format like a (natural) bridge
between the foreign image format and our simple data storage.


\begin{wideverbatim}

(setq *Ppm (ppmRead '("convert" "img.jpg" "ppm:-")))

\end{wideverbatim}

\pagebreak{}
\section*{Bitmap/Write a PPM file}

Using the data storage type defined \emph{Basic bitmap storage} for
raster images, write the image to a PPM file (binary P6 prefered).

(Read \href{http://en.wikipedia.org/wiki/Netpbm\_format}{the
  definition of PPM file} on Wikipedia.)


\begin{wideverbatim}

(de ppmWrite (Ppm File)
   (out File
      (prinl "P6")
      (prinl (length (car Ppm)) " " (length Ppm))
      (prinl 255)
      (for Y Ppm (for X Y (apply wr X))) ) )

\end{wideverbatim}

\pagebreak{}
\section*{Bitwise IO}

The aim of this task is to write functions (or create a class if your
language is Object Oriented and you prefer) for reading and writing
sequences of bits. While the output of a \texttt{asciiprint "STRING"} is
the ASCII byte sequence ``S'', ``T'', ``R'', ``I'', ``N'', ``G'', the
output of a ``print'' of the bits sequence 0101011101010 (13 bits) must
be 0101011101010; real I/O is performed always \emph{quantized} by byte
(avoiding endianness issues and relying on underlying buffering for
performance), therefore you must obtain as output the bytes 0101 0111
0101 0\textbf{000} (bold bits are padding bits), i.e. in hexadecimal 57
50.

As test, you can implement a \textbf{rough} (e.g. don't care about error
handling or other issues) compression/decompression program for ASCII
sequences of bytes, i.e. bytes for which the most significant bit is
always unused, so that you can write seven bits instead of eight (each 8
bytes of input, we write 7 bytes of output).

These bit oriented I/O functions can be used to implement compressors
and decompressors; e.g. Dynamic and Static Huffman encodings use
variable length bits sequences, while LZW (see \emph{LZW compression})
use fixed or variable \emph{words} nine (or more) bits long.

\begin{itemize}
\item
  Limits in the maximum number of bits that can be written/read in a
  single read/write operation are allowed.
\end{itemize}

\begin{itemize}
\item
  Errors handling is not mandatory
\end{itemize}


\begin{wideverbatim}

(de write7bitwise (Lst)
   (let (Bits 0  Byte)
      (for N Lst
         (if (=0 Bits)
            (setq Bits 7  Byte (* 2 N))
            (wr (| Byte (>> (dec 'Bits) N)))
            (setq Byte (>> (- Bits 8) N)) ) )
      (unless (=0 Bits)
         (wr Byte) ) ) )

(de read7bitwise ()
   (make
      (let (Bits 0  Byte)
         (while (rd 1)
            (let N @
               (link
                  (if (=0 Bits)
                     (>> (one Bits) N)
                     (| Byte (>> (inc 'Bits) N)) ) )
               (setq Byte (\& 127 (>> (- Bits 7) N))) ) )
         (when (= 7 Bits)
            (link Byte) ) ) ) )


(out 'a (write7bitwise (127 0 127 0 127 0 127 0 127)))
(hd 'a)
(in 'a (println (read7bitwise)))

(out 'a (write7bitwise (0 127 0 127 0 127 0 127 0)))
(hd 'a)
(in 'a (println (read7bitwise)))

(out 'a (write7bitwise (mapcar char (chop "STRING"))))
(hd 'a)
(println (mapcar char (in 'a (read7bitwise))))

Output:

00000000  FE 03 F8 0F E0 3F 80 FE                          .....?..
(127 0 127 0 127 0 127 0)
00000000  01 FC 07 F0 1F C0 7F 00                          .......
(0 127 0 127 0 127 0 127)
00000000  A7 52 94 99 D1 C0                                .R....
("S" "T" "R" "I" "N" "G")

\end{wideverbatim}

\pagebreak{}
\section*{Bitwise operations}


\textbf{Basic Data Operation}\\ This is a basic data operation. It
represents a fundamental action on a basic data type.

You may see other such operations in the \emph{Basic Data Operations}
category, or:

\textbf{Integer Operations} \\
\emph{Arithmetic} \textbar{} \emph{Comparison}

\textbf{Boolean Operations} \\ \textbf{Bitwise} \textbar{} \emph{Logical}

\textbf{String Operations} \\
\emph{Concatenation} \textbar{} 
\emph{Interpolation} \textbar{}
\emph{Matching}

\textbf{Memory Operations} \\
\emph{Pointers \& references}
\textbar{} \emph{Addresses}

Write a routine to perform a bitwise AND, OR, and XOR on two integers, a
bitwise NOT on the first integer, a left shift, right shift, right
arithmetic shift, left rotate, and right rotate. All shifts and rotates
should be done on the first integer with a shift/rotate amount of the
second integer. If any operation is not available in your language, note
it.



\begin{wideverbatim}

PicoLisp has no specific word size. Numbers grow to arbitrary length. Therefore,
bitwise NOT, logical (non-arithmetic) SHIFTs, and rotate operations do not make
sense.

Bitwise AND:

: (\& 6 3)
-> 2

: (\& 7 3 1)
-> 1

Bitwise AND-Test (tests if all bits in the first argument are set in the
following arguments):

: (bit? 1 2)
-> NIL

: (bit? 6 3)
-> NIL

: (bit? 6 15 255)
-> 6

Bitwise OR:

: (| 1 2)
-> 3

: (| 1 2 4 8)
-> 15

Bitwise XOR:

: (x| 2 7)
-> 5

: (x| 2 7 1)
-> 4

Shift (right with a positive count, left with a negative count):

: (>> 1 8)
-> 4

: (>> 3 16)
-> 2

: (>> -3 16)
-> 128

: (>> -1 -16)
-> -32

\end{wideverbatim}

\pagebreak{}
\section*{Boolean values}


Show how to represent the boolean states ``true'' and ``false'' in a
language. If other objects represent ``true'' or ``false'' in
conditionals, note it.

Cf.

\begin{itemize}
\item
  \emph{Logical operations}
\end{itemize}



\begin{wideverbatim}

Like in all Lisps, the symbol 'NIL' denotes "false", any other value "true".

Some functions return the symbol 'T' for "true" if no other useful (non-NIL)
value is available in the given context. Note that 'NIL' and 'T' are written in
uppercase letters (PicoLisp is case-sensitive).

\end{wideverbatim}

\pagebreak{}
\section*{Boxing the compass}

Avast me hearties!

There be many a \href{http://www.talklikeapirate.com/howto.html}{land
lubber} that knows
\href{http://oxforddictionaries.com/view/entry/m\_en\_gb0550020\#m\_en\_gb0550020}{naught}
of the pirate ways and gives direction by degree! They know not how to
\href{http://en.wikipedia.org/wiki/Boxing\_the\_compass}{box the
compass}!

\textbf{Task description}

\begin{enumerate}
\item
  Create a function that takes a heading in degrees and returns the
  correct 32-point compass heading.
\item
  Use the function to print and display a table of Index, Compass point,
  and Degree; rather like the corresponding columns from, the first
  table of the
  \href{http://en.wikipedia.org/wiki/Boxing\_the\_compass}{wikipedia
  article}, but use only the following 33 headings as input:
\end{enumerate}

\texttt{{[}0.0, 16.87, 16.88, 33.75, 50.62, 50.63, 67.5, 84.37, 84.38, 101.25, 118.12, 118.13, 135.0, 151.87, 151.88, 168.75, 185.62, 185.63, 202.5, 219.37, 219.38, 236.25, 253.12, 253.13, 270.0, 286.87, 286.88, 303.75, 320.62, 320.63, 337.5, 354.37, 354.38{]}}.
(They should give the same order of points but are spread throughout the
ranges of acceptance).

Notes;

\begin{itemize}
\item
  The headings and indices can be calculated from this pseudocode:
\end{itemize}

\begin{verbatim}
for i in 0..32 inclusive:
    heading = i * 11.25
    case i %3:
      if 1: heading += 5.62; break
      if 2: heading -= 5.62; break
    end
    index = ( i mod 32) + 1
\end{verbatim}

\begin{itemize}
\item The column of indices can be thought of as an enumeration of the
  thirty two cardinal points (see \emph{talk page})..
\end{itemize}


\begin{wideverbatim}

(scl 3)

(setq *Compass                      # Build lookup table
   (let H -16.875
      (mapcar
         '((Str)
            (cons
               (inc 'H 11.25)       # Heading in degrees
               (pack                # Compass point
                  (replace (chop Str)
                     "N" "north"
                     "E" "east"
                     "S" "south"
                     "W" "west"
                     "b" " by " ) ) ) )
         '("N" "NbE" "N-NE" "NEbN" "NE" "NEbE" "E-NE" "EbN"
            "E" "EbS" "E-SE" "SEbE" "SE" "SEbS" "S-SE" "SbE"
            "S" "SbW" "S-SW" "SWbS" "SW" "SWbW" "W-SW" "WbS"
            "W" "WbN" "W-NW" "NWbW" "NW" "NWbN" "N-NW" "NbW"
            "N" ) ) ) )

(de heading (Deg)
   (rank (\% Deg 360.00) *Compass) )

(for I (range 0 32)
   (let H (* I 11.25)
      (case (\% I 3)
         (1 (inc 'H 5.62))
         (2 (dec 'H 5.62)) )
      (tab (3 1 -18 8)
         (inc (\% I 32))
         NIL
         (cdr (heading H))
         (round H 2) ) ) )

\end{wideverbatim}

\begin{wideverbatim}

Output:

  1 north                 0.00
  2 north by east        16.87
  3 north-northeast      16.88
  4 northeast by north   33.75
  5 northeast            50.62
  6 northeast by east    50.63
  7 east-northeast       67.50
  8 east by north        84.37
  9 east                 84.38
 10 east by south       101.25
 11 east-southeast      118.12
 12 southeast by east   118.13
 13 southeast           135.00
 14 southeast by south  151.87
 15 south-southeast     151.88
 16 south by east       168.75
 17 south               185.62
 18 south by west       185.63
 19 south-southwest     202.50
 20 southwest by south  219.37
 21 southwest           219.38
 22 southwest by west   236.25
 23 west-southwest      253.12
 24 west by south       253.13
 25 west                270.00
 26 west by north       286.87
 27 west-northwest      286.88
 28 northwest by west   303.75
 29 northwest           320.62
 30 northwest by north  320.63
 31 north-northwest     337.50
 32 north by west       354.37
  1 north               354.38

\end{wideverbatim}

\pagebreak{}
\section*{Break OO privacy}

Show how to access private or protected members of a class in an object
oriented language from outside an instance of the class, without calling
non-private or non-protected members of the class as a proxy.

\emph{Note that this is almost universally regarded as unidiomatic at
best, and poor programming practice at worst.}



\begin{wideverbatim}

PicoLisp uses [http://software-lab.de/doc/ref.html#transient "transient
symbols"] for variables, functions, methods etc. inaccessible from other parts
of the program. Lexically, a transient symbol is enclosed by double quotes.

The only way to access a transient symbol outside its namespace is to search for
its name in other (public) structures. This is done by the
'[http://software-lab.de/doc/refL.html#loc loc]' function.

(class +Example)
# "_name"

(dm T (Name)
   (=: "_name" Name) )

(dm string> ()
   (pack "Hello, I am " (: "_name")) )

(====)  # Close transient scope

(setq Foo (new '(+Example) "Eric"))

\end{wideverbatim}

\begin{wideverbatim}
Test:

: (string> Foo)                        # Access via method call
-> "Hello, I am Eric"

: (get Foo '"_name")                   # Direct access doesn't work
-> NIL

: (get Foo (loc "_name" +Example))     # Locating the transient symbol works
-> "Eric"

: (put Foo (loc "_name" +Example) "Edith")
-> "Edith"

: (string> Foo)                        # Ditto
-> "Hello, I am Edith"

: (get Foo '"_name")
-> NIL

: (get Foo (loc "_name" +Example))
-> "Edith"

\end{wideverbatim}

\pagebreak{}
\section*{Brownian tree}


Generate and draw a
\href{http://en.wikipedia.org/wiki/Brownian\_tree}{Brownian Tree}.

A Brownian Tree is generated as a result of an initial seed, followed by
the interaction of two processes.

\begin{enumerate}
\item
  The initial ``seed'' is placed somewhere within the field. Where is
  not particularly important; it could be randomized, or it could be a
  fixed point.
\item
  Particles are injected into the field, and are individually given a
  (typically random) motion pattern.
\item
  When a particle collides with the seed or tree, its position is fixed,
  and it's considered to be part of the tree.
\end{enumerate}

Because of the lax rules governing the random nature of the particle's
placement and motion, no two resulting trees are really expected to be
the same, or even necessarily have the same general shape.



\begin{wideverbatim}

(load "@lib/simul.l")

(de brownianTree (File Size Cnt)
   (let Img (grid Size Size)
      (put Img (/ Size 2) (/ Size 2) 'pix T)
      (use (P Q)
         (do Cnt
            (setq P (get Img (rand 1 Size) (rand 1 Size)))
            (loop
               (setq Q ((if2 (rand T) (rand T) north east south west) P))
               (T (; Q pix) (put P 'pix T))
               (setq P (or Q (get Img (rand 1 Size) (rand 1 Size)))) ) ) )
      (out "img.pbm"
         (prinl "P1")
         (prinl Size " " Size)
         (for L Img
            (for This L
               (prin (if (: pix) 1 0)) )
            (prinl) ) ) ) )

Use:

(brownianTree "img.pbm" 300 9000)
(call 'display "img.pbm")

\end{wideverbatim}

% \pagebreak{}
% \section*{Bulls and cows}

% \href{http://en.wikipedia.org/wiki/Bulls\_and\_Cows}{This} is an old
% game played with pencil and paper that was later implemented on
% computer.

% The task is for the program to create a four digit random number from
% the digits 1 to 9, without duplication. The program should ask for
% guesses to this number, reject guesses that are malformed, then print
% the score for the guess.

% The score is computed as:

% \begin{enumerate}
% \item
%   The player wins if the guess is the same as the randomly chosen
%   number, and the program ends.
% \item
%   A score of one \textbf{bull} is accumulated for each digit in the
%   guess that equals the corresponding digit in the randomly chosen
%   initial number.
% \item
%   A score of one \textbf{cow} is accumulated for each digit in the guess
%   that also appears in the randomly chosen number, but in the wrong
%   position.
% \end{enumerate}

% Cf,

% \begin{itemize}
% \item
%   \emph{Bulls and cows/Player}


%   The task is to write a \emph{player} of the \emph{Bulls and Cows
%     game}, rather than a scorer. The player should give intermediate
%   answers that respect the scores to previous attempts.

% One method is to generate a list of all possible numbers that could be
% the answer, then to prune the list by keeping only those numbers that
% would give an equivalent score to how your last guess was scored. Your
% next guess can be any number from the pruned list.\\ Either you guess
% correctly or run out of numbers to guess, which indicates a problem with
% the scoring.

% Cf,

% \begin{itemize}
% \item
%   \emph{Bulls and cows}
% \item
%   \emph{Guess the number}
% \item
%   \emph{Guess the number/With Feedback (Player)}
% \end{itemize}


% \item
%   \emph{Guess the number}
% \item
%   \emph{Guess the number/With Feedback}
% \end{itemize}


% \begin{wideverbatim}

% (de ok? (N)
%    (let D
%       (mapcar 'format (chop N))
%       (and
%          (num? N)
%          (not (member 0 D))
%          (= 4 (length D))
%          (= D (uniq D))
%          D ) ) )

% (de init-cows ()
%    (until (setq *Hidden (ok? (rand 1234 9876)))) )

% (de guess (N)
%    (let D (ok? N)
%       (if D
%          (let Bulls (cnt '= D *Hidden)
%             (if (= 4 Bulls)
%                " You guessed it!"
%                (let Cows (- (cnt '((N) (member N *Hidden)) D ) Bulls )
%                   (pack Bulls " bulls, " Cows " cows") ) ) )
%          " Bad guess! (4 unique digits, 1-9)" ) ) )

% \end{wideverbatim}


\pagebreak{}
\section*{Bulls and cows/Player}


The task is to write a \emph{player} of the
\emph{Bulls and Cows game}, rather than a
scorer. The player should give intermediate answers that respect the
scores to previous attempts.

One method is to generate a list of all possible numbers that could be
the answer, then to prune the list by keeping only those numbers that
would give an equivalent score to how your last guess was scored. Your
next guess can be any number from the pruned list.\\ Either you guess
correctly or run out of numbers to guess, which indicates a problem with
the scoring.

Cf,

\begin{itemize}
\item
  \emph{Bulls and cows}
\item
  \emph{Guess the number}
\item
  \emph{Guess the number/With Feedback (Player)}
\end{itemize}

\begin{wideverbatim}

(load "@lib/simul.l")

(de bullsAndCows ()
   (let Choices (shuffle (mapcan permute (subsets 4 (range 1 9))))
      (use (Guess Bulls Cows)
         (loop
            (prinl "Guessing " (setq Guess (pop 'Choices)))
            (prin "How many bulls and cows? ")
            (setq Bulls (read)  Cows (read))
            (setq Choices
               (filter
                  '((C)
                     (let B (cnt = Guess C)
                        (and
                           (= Bulls B)
                           (= Cows (- (length (sect Guess C)) B)) ) ) )
                  Choices ) )
            (NIL Choices "No matching solution")
            (NIL (cdr Choices) (pack "The answer is " (car Choices))) ) ) ) )

\end{wideverbatim}

\begin{wideverbatim}

Output:

: (bullsAndCows)
Guessing 4217
How many bulls and cows? 0 2
Guessing 5762
How many bulls and cows? 1 1
Guessing 9372
How many bulls and cows? 0 1
Guessing 7864
How many bulls and cows? 1 2
Guessing 8754
How many bulls and cows? 0 2
-> "The answer is 2468"

\end{wideverbatim}



% put{referenc}
